\documentclass[conference]{IEEEtran}
\IEEEoverridecommandlockouts

\makeatletter
\newcommand{\newlineauthors}{%
  \end{@IEEEauthorhalign}\hfill\mbox{}\par
  \mbox{}\hfill\begin{@IEEEauthorhalign}
}
\makeatother

\usepackage{cite}
\usepackage{amsmath,amssymb,amsfonts}
\usepackage{algorithmic}
\usepackage{algorithm}
\usepackage{graphicx}
\usepackage{textcomp}
\usepackage{xcolor}
\usepackage{booktabs}
\usepackage{multirow}
\usepackage[colorlinks=true, linkcolor=black, citecolor=black, urlcolor=black]{hyperref}

\def\BibTeX{{\rm B\kern-.05em{\sc i\kern-.025em b}\kern-.08em
    T\kern-.1667em\lower.7ex\hbox{E}\kern-.125emX}}

\begin{document}

\title{Cache-Aided NOMA for 6G Networks: Deep Reinforcement Learning for\\
CIC-Aware Proactive Cache Placement}

\author{
\IEEEauthorblockN{Akash Mandal}
\IEEEauthorblockA{\textit{Dept. of Information Technology}\\
\textit{Techno Main Salt Lake}\\
Kolkata, India\\
akash.lawrencian2020@gmail.com}
\and
\IEEEauthorblockN{Arpita Guchhait}
\IEEEauthorblockA{\textit{Dept. of Information Technology}\\
\textit{Techno Main Salt Lake}\\
Kolkata, India\\
arpitaguchhait2020@gmail.com}
\and
\IEEEauthorblockN{Manojit Pal}
\IEEEauthorblockA{\textit{Dept. of Information Technology}\\
\textit{Techno Main Salt Lake}\\
Kolkata, India\\
manojitmunmun@gmail.com}
\newlineauthors
\IEEEauthorblockN{Vaibhav Anand}
\IEEEauthorblockA{\textit{Dept. of Information Technology}\\
\textit{Techno Main Salt Lake}\\
Kolkata, India\\
anandvaibhav0107@gmail.com}
\and
\IEEEauthorblockN{Dr. Biswajit Ghosh}
\IEEEauthorblockA{\textit{Dept. of Information Technology}\\
\textit{Techno Main Salt Lake}\\
Kolkata, India \\
biswanerist@gmail.com}
}

\maketitle

%=============================================================
% ABSTRACT
%=============================================================
\begin{abstract}
This paper proposes a Deep Reinforcement Learning (DRL) based proactive caching framework for Cache-aided Non-Orthogonal Multiple Access (NOMA) downlink systems targeting 6G networks. While caching and NOMA are promising complementary technologies, traditional reactive caching policies fail to optimize for Cache-aided Interference Cancellation (CIC) opportunities. The proposed scheme employs a Double Dueling Deep Q-Network (DQN) agent with a novel binary action space to make online cache placement decisions. By incorporating a CIC-aware reward function, the agent is explicitly incentivized to proactively cache files that enable interference suppression for paired NOMA users. Simulated under a realistic mixed fading channel model with imperfect Successive Interference Cancellation (SIC), results over 5 million training steps and 100 Monte Carlo evaluation runs demonstrate exceptional performance. The proposed DQN scheme achieves a cache hit rate of 70.38\%, a CIC benefit rate of 3.78 (a 5.28x improvement over the best dynamic baseline), and a sum-rate of 18.17 bps/Hz at an SNR of 30 dB. Furthermore, the proposed policy reduces system outage probability by 14x compared to the no-caching baseline, performing within 0.54\% of the oracle-optimal Top-K policy without relying on prior knowledge of the content popularity distribution.
\end{abstract}

\begin{IEEEkeywords}
cache-aided NOMA, deep reinforcement learning, 6G, proactive caching, CIC, Double DQN, interference cancellation, edge caching
\end{IEEEkeywords}


%=============================================================
% I. INTRODUCTION
%=============================================================
\section{Introduction}

Next-generation (6G) wireless networks are expected to meet unprecedented demands for ultra-low latency, massive connectivity, and high spectral efficiency \cite{b_ding_survey}. Two complementary technologies---Non-Orthogonal Multiple Access (NOMA) and edge caching---have independently demonstrated significant gains in spectral efficiency and backhaul reduction, respectively. NOMA enables multiple users to share the same time-frequency resources via power-domain multiplexing and Successive Interference Cancellation (SIC) \cite{b_dai_noma, b_ding_survey}, while edge caching reduces access latency and backhaul load by proactively placing popular content near the users \cite{b_maddah_ali, b_golrezaei}.

The joint paradigm---\textit{Cache-Aided NOMA}---has emerged as a highly promising architecture for 6G, where cached content can further enable Cache-aided Interference Cancellation (CIC). When a file requested by a paired NOMA user is locally available in the cache, the multi-user interference term from that user can be perfectly reconstructed and canceled without requiring error-prone SIC decoding. This fundamentally reduces outage probability and improves system-level fairness \cite{b_survey_bepari, b_cao_joint}.

However, the full potential of CIC depends critically on the \textit{quality} of the cached content---specifically, whether the cached files are the ones likely to be requested by paired NOMA users in the near future. Static and reactive policies (e.g., Top-K, Least Recently Used (LRU), Least Frequently Used (LFU)) do not account for this cross-layer NOMA coupling, resulting in severely suboptimal CIC exploitation. Deep Reinforcement Learning (DRL) offers an appealing solution: an intelligent agent can learn, through continuous environmental interaction, to dynamically cache files that simultaneously maximize the cache hit rate and CIC opportunities.

\subsection*{Contributions}
The primary contributions of this paper are summarized as follows:
\begin{itemize}
    \item We design a novel NOMA-aware state representation (14-dimensional) for the DRL agent that jointly encodes content popularity, cache occupancy, continuous channel statistics, and recent NOMA pairing history.
    \item We present a binary action space redesign that reduces the traditional DQN eviction action space from $\mathcal{O}(N_{\text{files}})$ to $\mathcal{O}(1)$, making DRL remarkably tractable for large, real-world content catalogs.
    \item We propose a specialized CIC-incentivized immediate reward function that explicitly rewards cache placement decisions enabling interference cancellation for paired NOMA users.
    \item We conduct a comprehensive simulation study under realistic mixed fading channels with imperfect SIC, comparing the DQN scheme against four baselines (Top-K Zipf, LRU, LFU, Random) and a No-Cache reference over a full signal-to-noise ratio (SNR) sweep from $-10$ to $30$ dB.
\end{itemize}


%=============================================================
% II. RELATED WORK
%=============================================================
\section{Related Work}

Edge caching for wireless systems has been extensively studied since early empirical observations demonstrated that content popularity follows Zipf-like distributions \cite{b_breslau_zipf}. Fundamental limits of caching---formalized in the coded-caching framework---proved that coordinated placement enables multicast opportunities and global caching gains \cite{b_maddah_ali}. Architectures such as Femtocaching and device-to-device (D2D) caching have subsequently been proposed to exploit spatial locality to reduce core network load \cite{b_golrezaei, b_ji_d2d}.

The convergence of edge caching and NOMA has recently been analyzed from stochastic geometry \cite{b_zhang_stochastic} and joint optimization perspectives \cite{b_cao_joint, b_guo_edge}. A comprehensive recent survey \cite{b_survey_bepari} identified CIC as a key unexplored dimension in cache-aided NOMA. While cached content can eliminate an interferer's signal at the receiver without SIC, this cross-layer benefit has not been fully exploited in conventional machine-learning-based placement policies.

Machine-learning approaches---ranging from statistical predictors to deep-RL methods---have been increasingly proposed for adaptive cache storage \cite{b_chen_ml, b_elghariani_pred}. Li \textit{et al.} \cite{b_li_drl} studied DRL-based dynamic caching in NOMA networks but did not incorporate a CIC-aware reward feedback mechanism or the binary action-space abstraction necessary for scalable catalog deployments. Cooperative and UAV-relay assisted NOMA designs have further been studied for cell-edge coverage \cite{b_bepari_partial, b_liu_uav}. Our work explicitly differs by formulating the first CIC-incentivized DRL proactive caching policy with direct NOMA-layer state feedback mapping.


%=============================================================
% III. SYSTEM MODEL
%=============================================================
%====================================================
\section{System Model}
%====================================================

%====================================================
\subsection{Channel Model}
%====================================================

The wireless channel between the BS and each user is assumed to be frequency-flat within an individual Resource Block (RB). Let $h_{k,m}$ denote the complex channel coefficient between the BS and user $k$ over the $m$-th RB. The overall channel consists of three major propagation components, namely distance-dependent path loss, log-normal shadowing, and small-scale fading. Therefore, the instantaneous channel gain can be expressed as

\begin{equation}
|h_{k,m}|^2=
\frac{\xi_k |g_{k,m}|^2}
{PL_k},
\label{eq:channel_gain}
\end{equation}

where

\begin{itemize}
\item $PL_k$ denotes the large-scale path loss,
\item $\xi_k$ represents the shadowing coefficient,
\item $g_{k,m}$ denotes the small-scale fading coefficient.
\end{itemize}

The distance-dependent path loss is modeled as

\begin{equation}
PL_k=d_k^{\beta},
\label{eq:pathloss}
\end{equation}

where $d_k$ is the distance between the BS and user $k$, while $\beta$ is the path-loss exponent.

The shadowing component follows a log-normal distribution

\begin{equation}
10\log_{10}(\xi_k)\sim\mathcal{N}(0,\sigma_{sh}^2),
\end{equation}

where $\sigma_{sh}$ denotes the shadowing standard deviation.

The small-scale fading coefficient is assumed to follow Rayleigh fading, i.e.,

\begin{equation}
g_{k,m}\sim\mathcal{CN}(0,1),
\end{equation}

which results in exponentially distributed channel power gains.

The received signal at user $k$ over the $m$-th RB is

\begin{equation}
y_{k,m}
=
h_{k,m}x_m+n_{k,m},
\label{eq:received_signal}
\end{equation}

where $x_m$ denotes the transmitted superposition signal and

\begin{equation}
n_{k,m}\sim\mathcal{CN}(0,\sigma^2)
\end{equation}

is the additive white Gaussian noise (AWGN).

To facilitate NOMA transmission, users are arranged in ascending order according to their instantaneous channel gains

\begin{equation}
|h_{1,m}|^2
\le
|h_{2,m}|^2
\le
\cdots
\le
|h_{K,m}|^2.
\label{eq:ordering}
\end{equation}

For each RB, the user having the minimum channel gain is referred to as the weak user ($w$), whereas the user with the highest channel gain is referred to as the strong user ($s$). The proposed system adopts an extreme pairing strategy in which the strongest user is paired with the weakest user to maximize the channel gain difference and improve the effectiveness of power-domain multiplexing.

%====================================================
\subsection{Cache Placement Model}
%====================================================

The proposed framework adopts a proactive cache placement strategy in which frequently requested contents are stored at user devices during off-peak network hours. The cache placement phase is performed before the actual transmission process and therefore does not consume radio resources during data delivery.

Let the entire content library be represented by

\begin{equation}
\mathcal{F}
=
\{f_1,f_2,\cdots,f_N\},
\end{equation}

where each file has identical size.

Each user is equipped with a local cache memory of capacity

\begin{equation}
C<N,
\end{equation}

allowing only a subset of the library to be stored locally.

The popularity of the files follows the Zipf distribution

\begin{equation}
p_n=
\frac{n^{-\alpha}}
{\sum_{i=1}^{N}i^{-\alpha}},
\qquad
n=1,\cdots,N,
\label{eq:zipf_distribution}
\end{equation}

where $\alpha$ controls the popularity skewness.

During the cache placement phase, the BS stores the most popular contents at users according to the cache capacity constraint

\begin{equation}
\sum_{n=1}^{N}c_{k,n}\le C,
\label{eq:cache_constraint}
\end{equation}

where

\begin{equation}
c_{k,n}
=
\begin{cases}
1, & \text{if file }f_n\text{ is cached at user }k,\\
0, & \text{otherwise}.
\end{cases}
\label{eq:cache_indicator}
\end{equation}

When a user requests a file already available in its cache, the requested content is served locally without occupying wireless resources. Otherwise, the BS delivers the requested file through the downlink NOMA transmission.

Unlike conventional cache-enabled systems, the proposed Hybrid SIC/CIC framework also exploits cached information to suppress multi-user interference. Consider a NOMA pair consisting of a weak user ($w$) and a strong user ($s$). Let

\begin{equation}
c_w=
\begin{cases}
1,& \text{if the strong user's requested file is cached at the weak user},\\
0,& \text{otherwise},
\end{cases}
\end{equation}

and

\begin{equation}
c_s=
\begin{cases}
1,& \text{if the weak user's requested file is cached at the strong user},\\
0,& \text{otherwise}.
\end{cases}
\end{equation}

The binary variables $c_w$ and $c_s$ determine whether Cache-Aided Interference Cancellation (CIC) can be activated at the corresponding receiver. If $c_w=1$, the weak user can reconstruct and perfectly remove the strong user's interference directly from its local cache. Similarly, when $c_s=1$, the strong user can eliminate the weak user's signal without relying on successive interference cancellation.

Accordingly, four different transmission scenarios are possible:

\begin{enumerate}
\item $c_w=0,\;c_s=0$: Conventional NOMA with SIC.
\item $c_w=1,\;c_s=0$: Weak-user Cache-Aided Interference Cancellation.
\item $c_w=0,\;c_s=1$: Strong-user Cache-Aided Interference Cancellation.
\item $c_w=1,\;c_s=1$: Full Hybrid CIC, where both users cancel interference using cached information.
\end{enumerate}

This cache-aware transmission strategy enables the proposed system to dynamically switch between conventional SIC and CIC according to the cache status of each NOMA pair, thereby improving decoding reliability, reducing residual interference, and lowering the overall outage probability.

%====================================================
\subsection{User Pairing Strategy}
%====================================================

To improve spectral efficiency and exploit the power-domain multiplexing capability of NOMA, two users are grouped together on each Resource Block (RB). The proposed framework adopts an extreme user pairing strategy, where the user experiencing the strongest channel condition is paired with the user having the weakest channel condition within the same RB. Such pairing maximizes the channel gain difference between paired users, thereby improving power allocation efficiency and decoding performance.

Let the ordered channel gains over the $m$-th RB be

\begin{equation}
|h_{1,m}|^2
\le
|h_{2,m}|^2
\le
\cdots
\le
|h_{K,m}|^2.
\end{equation}

The paired users are denoted as

\begin{equation}
(w,s),
\end{equation}

where

\begin{equation}
|h_{w,m}|^2<|h_{s,m}|^2.
\end{equation}

The weak user is allocated a larger transmit power, while the strong user receives a smaller power allocation. Let

\begin{equation}
\alpha_{w,m}+\alpha_{s,m}=1,
\end{equation}

where

\begin{equation}
\alpha_{w,m}>\alpha_{s,m}.
\end{equation}

This power allocation guarantees that the weak user's Quality-of-Service (QoS) requirement is satisfied while allowing the strong user to decode the weak user's signal using successive interference cancellation (SIC) whenever cache-assisted cancellation is unavailable.

Each Resource Block serves only one NOMA pair. The RB allocation indicator is defined as

\begin{equation}
\beta_{k,m}
=
\begin{cases}
1,&\text{if user }k\text{ is assigned to RB }m,\\
0,&\text{otherwise}.
\end{cases}
\end{equation}

The RB assignment satisfies

\begin{equation}
\sum_{k=1}^{K}\beta_{k,m}\le2,
\qquad
\forall m,
\end{equation}

ensuring that at most two users share the same RB.

%====================================================
\subsection{Cache-Aided NOMA Transmission}
%====================================================

After user pairing, the BS simultaneously transmits the requested information to both users through power-domain superposition coding.

The transmitted signal over the $m$-th RB is

\begin{equation}
x_m
=
\sqrt{\alpha_{w,m}P_m}x_w
+
\sqrt{\alpha_{s,m}P_m}x_s,
\label{eq:superposition}
\end{equation}

where

\begin{itemize}
\item $P_m$ denotes the transmit power allocated to the $m$-th RB,
\item $x_w$ and $x_s$ represent the unit-power information symbols intended for the weak and strong users,
\item $\alpha_{w,m}$ and $\alpha_{s,m}$ denote the corresponding power allocation coefficients.
\end{itemize}

The received signal at user $k$ is

\begin{equation}
y_{k,m}
=
h_{k,m}
\left(
\sqrt{\alpha_{w,m}P_m}x_w
+
\sqrt{\alpha_{s,m}P_m}x_s
\right)
+n_{k,m},
\end{equation}

where

\begin{equation}
n_{k,m}\sim\mathcal{CN}(0,\sigma^2).
\end{equation}

Unlike conventional NOMA systems, the proposed framework utilizes cached contents not only to reduce backhaul traffic but also to suppress multi-user interference before signal detection.

Depending on the cache status of the paired users, the receiver dynamically switches between conventional SIC and Cache-Aided Interference Cancellation (CIC). Consequently, the received signal can be processed under four different operating conditions, thereby improving the overall decoding reliability and reducing transmission outages.

%====================================================
\subsection{Hybrid SIC/CIC Receiver Model}
%====================================================

In conventional NOMA systems, the weak user directly decodes its own signal by treating the strong user's transmission as interference. The strong user first decodes the weak user's signal and subsequently removes it using successive interference cancellation (SIC). However, practical SIC suffers from imperfect signal reconstruction, which leaves residual interference and degrades system performance.

To overcome this limitation, the proposed framework introduces a Hybrid SIC/CIC receiver that dynamically selects the interference cancellation mechanism according to the cache status of each user.

Let

\begin{equation}
c_w=
\begin{cases}
1,&\text{if the strong user's requested file is cached at the weak user},\\
0,&\text{otherwise},
\end{cases}
\end{equation}

and

\begin{equation}
c_s=
\begin{cases}
1,&\text{if the weak user's requested file is cached at the strong user},\\
0,&\text{otherwise}.
\end{cases}
\end{equation}

The variables $c_w$ and $c_s$ determine whether cache-assisted interference cancellation can be activated.

Accordingly, the receiver operates under one of the following four transmission modes.

\subsubsection*{Mode 1: Conventional SIC}

When

\[
c_w=0,\qquad c_s=0,
\]

no useful cached information is available. Therefore, both users follow the conventional NOMA decoding procedure. The weak user treats the strong-user signal as interference, whereas the strong user performs imperfect SIC.

\subsubsection*{Mode 2: Weak-User CIC}

When

\[
c_w=1,\qquad c_s=0,
\]

the weak user reconstructs the interfering signal directly from its local cache and removes it before decoding its desired signal. The strong user continues to employ conventional SIC.

\subsubsection*{Mode 3: Strong-User CIC}

When

\[
c_w=0,\qquad c_s=1,
\]

the strong user already possesses the weak user's requested file in its cache. Therefore, it reconstructs the interference perfectly without performing SIC. The weak user still treats the strong-user transmission as interference.

\subsubsection*{Mode 4: Full Hybrid CIC}

When

\[
c_w=1,\qquad c_s=1,
\]

both users have access to each other's requested content through local caching. Consequently, each receiver reconstructs and removes the interfering signal independently, eliminating the need for SIC completely.

The proposed Hybrid SIC/CIC receiver therefore provides a generalized interference management framework that seamlessly switches between conventional SIC and cache-assisted interference cancellation according to the instantaneous cache availability. This adaptive receiver improves decoding reliability, reduces residual interference, and forms the basis for the unified SINR and outage analysis presented in the subsequent section.
```latex
%====================================================
\subsection{Unified SINR Analysis}
%====================================================

The proposed Hybrid SIC/CIC receiver operates according to the cache availability of the paired users. Instead of deriving separate SINR expressions for each operating mode, a unified mathematical model is developed to represent all possible transmission scenarios.

Let

\begin{equation}
c_w,c_s\in\{0,1\},
\end{equation}

be binary cache indicators, where

\begin{itemize}
\item $c_w=1$ indicates that the weak user stores the strong user's requested content,
\item $c_s=1$ indicates that the strong user stores the weak user's requested content.
\end{itemize}

If the required interfering signal is available in the local cache, the corresponding interference can be perfectly reconstructed and removed before signal detection.

Therefore, the received SINR of the weak user can be expressed in a unified form as

\begin{equation}
\gamma_{w,m}
=
\frac{\alpha_{w,m}P_m|h_{w,m}|^2}
{(1-c_w)\alpha_{s,m}P_m|h_{w,m}|^2+\sigma^2}.
\label{eq:UnifiedWeakSINR}
\end{equation}

Similarly, the SINR of the strong user is

\begin{equation}
\gamma_{s,m}
=
\frac{\alpha_{s,m}P_m|h_{s,m}|^2}
{(1-c_s)\zeta\alpha_{w,m}P_m|h_{s,m}|^2+\sigma^2},
\label{eq:UnifiedStrongSINR}
\end{equation}

where $\zeta$ denotes the residual interference factor caused by imperfect SIC.

The above equations automatically represent all four receiver operating modes.

\textbf{Case 1: Conventional SIC}

When

\[
c_w=0,\qquad c_s=0,
\]

(17) and (18) become

\begin{equation}
\gamma_w=
\frac{\alpha_wP|h_w|^2}
{\alpha_sP|h_w|^2+\sigma^2},
\end{equation}

\begin{equation}
\gamma_s=
\frac{\alpha_sP|h_s|^2}
{\zeta\alpha_wP|h_s|^2+\sigma^2},
\end{equation}

which correspond to the conventional imperfect SIC receiver.

\textbf{Case 2: Weak-User CIC}

For

\[
c_w=1,\qquad c_s=0,
\]

the interference observed by the weak user is completely removed,

\begin{equation}
\gamma_w^{CIC}
=
\frac{\alpha_wP|h_w|^2}
{\sigma^2},
\end{equation}

whereas the strong user still performs conventional SIC.

\textbf{Case 3: Strong-User CIC}

When

\[
c_w=0,\qquad c_s=1,
\]

the strong user performs perfect interference cancellation using its cached content,

\begin{equation}
\gamma_s^{CIC}
=
\frac{\alpha_sP|h_s|^2}
{\sigma^2},
\end{equation}

while the weak user follows the conventional NOMA decoding procedure.

\textbf{Case 4: Full Hybrid CIC}

If

\[
c_w=1,\qquad c_s=1,
\]

both users eliminate interference using cached information and their SINRs become

\begin{equation}
\gamma_w
=
\frac{\alpha_wP|h_w|^2}
{\sigma^2},
\end{equation}

\begin{equation}
\gamma_s
=
\frac{\alpha_sP|h_s|^2}
{\sigma^2}.
\end{equation}

Hence, the proposed unified formulation provides a generalized SINR model for both conventional SIC and cache-assisted interference cancellation.

%====================================================
\subsection{Achievable Rate Analysis}
%====================================================

Based on Shannon's channel capacity theorem, the instantaneous achievable data rate of user $k$ over the $m$-th Resource Block is

\begin{equation}
R_{k,m}
=
B
\log_2
\left(
1+\gamma_{k,m}
\right),
\end{equation}

where $B$ denotes the bandwidth allocated to one RB.

Accordingly, the achievable rates of the weak and strong users are

\begin{equation}
R_w
=
B
\log_2
(1+\gamma_w),
\end{equation}

and

\begin{equation}
R_s
=
B
\log_2
(1+\gamma_s),
\end{equation}

respectively.

The achievable sum rate of one NOMA pair is

\begin{equation}
R_{pair}
=
R_w+R_s.
\end{equation}

For a network containing $L$ independent NOMA pairs, the overall system throughput is

\begin{equation}
R_{sum}
=
\sum_{i=1}^{L}
R_{pair,i}.
\end{equation}

Since the proposed Hybrid SIC/CIC receiver suppresses multi-user interference whenever cached information is available, both $\gamma_w$ and $\gamma_s$ increase. Consequently, the achievable data rate also improves compared with the conventional SIC-based receiver.

%====================================================
\subsection{Weak User Outage Probability}
%====================================================

Reliable communication is achieved only when the received SINR exceeds a predefined threshold. Let the minimum Quality-of-Service (QoS) requirement be denoted by the target transmission rate $R_{th}$.

The corresponding SINR threshold is

\begin{equation}
\gamma_{th}
=
2^{R_{th}}-1.
\end{equation}

An outage event occurs whenever the received SINR falls below this threshold.

Therefore, the outage probability of the weak user is defined as

\begin{equation}
P_{out,w}
=
Pr
\left(
\gamma_w
<
\gamma_{th}
\right).
\end{equation}

Substituting the unified SINR expression from (17),

\begin{equation}
P_{out,w}
=
Pr
\left(
\frac{\alpha_wP|h_w|^2}
{(1-c_w)\alpha_sP|h_w|^2+\sigma^2}
<
\gamma_{th}
\right).
\end{equation}

For the conventional SIC receiver,

\[
c_w=0,
\]

the outage probability becomes

\begin{equation}
P_{out,w}^{SIC}
=
Pr
\left(
\frac{\alpha_wP|h_w|^2}
{\alpha_sP|h_w|^2+\sigma^2}
<
\gamma_{th}
\right).
\end{equation}

When cache-aided interference cancellation is activated,

\[
c_w=1,
\]

the interference term disappears and the outage probability reduces to

\begin{equation}
P_{out,w}^{CIC}
=
Pr
\left(
\frac{\alpha_wP|h_w|^2}
{\sigma^2}
<
\gamma_{th}
\right).
\end{equation}

Since the denominator contains only thermal noise, the weak user experiences a considerably lower outage probability under the proposed Hybrid SIC/CIC receiver than under the conventional SIC-based NOMA system.
```
```latex
%====================================================
\subsection{Strong User Outage Probability}
%====================================================

The strong user first removes the weak user's signal before decoding its own information. In the conventional NOMA receiver, this process is performed using Successive Interference Cancellation (SIC). However, practical SIC is not perfect and leaves a small amount of residual interference due to channel estimation errors and imperfect signal reconstruction.

Considering imperfect SIC, the outage probability of the strong user is defined as

\begin{equation}
P_{out,s}
=
\Pr
\left(
\gamma_s
<
\gamma_{th}
\right),
\end{equation}

where $\gamma_{th}=2^{R_{th}}-1$ denotes the target SINR threshold.

Substituting the unified SINR expression in (\ref{eq:UnifiedStrongSINR}), the outage probability becomes

\begin{equation}
P_{out,s}
=
\Pr
\left(
\frac{\alpha_sP|h_s|^2}
{(1-c_s)\zeta\alpha_wP|h_s|^2+\sigma^2}
<
\gamma_{th}
\right).
\end{equation}

Depending on the cache status of the strong user, two different decoding cases are obtained.

\subsubsection*{Case 1: Conventional SIC}

When the strong user does not store the weak user's requested content, i.e.,

\[
c_s=0,
\]

the receiver employs imperfect SIC and the outage probability is

\begin{equation}
P_{out,s}^{SIC}
=
\Pr
\left(
\frac{\alpha_sP|h_s|^2}
{\zeta\alpha_wP|h_s|^2+\sigma^2}
<
\gamma_{th}
\right).
\end{equation}

The residual interference coefficient $\zeta$ limits the achievable SINR and therefore increases the outage probability.

\subsubsection*{Case 2: Cache-Aided Interference Cancellation}

When the strong user already stores the weak user's requested content,

\[
c_s=1,
\]

the interfering signal is reconstructed directly from the cache without performing SIC. Consequently, the residual interference is completely eliminated and the received SINR becomes

\begin{equation}
\gamma_s^{CIC}
=
\frac{\alpha_sP|h_s|^2}
{\sigma^2}.
\end{equation}

The corresponding outage probability is

\begin{equation}
P_{out,s}^{CIC}
=
\Pr
\left(
\frac{\alpha_sP|h_s|^2}
{\sigma^2}
<
\gamma_{th}
\right).
\end{equation}

Since the residual interference term is completely removed, the proposed Hybrid SIC/CIC receiver significantly improves the decoding reliability of the strong user compared with the conventional SIC-based receiver.

%====================================================
\subsection{Pair Outage Probability}
%====================================================

A successful NOMA transmission requires that both the weak and the strong users correctly decode their intended signals. Therefore, a transmission outage occurs whenever at least one user experiences an outage.

Let

\begin{equation}
P_{out,w}
\end{equation}

and

\begin{equation}
P_{out,s}
\end{equation}

denote the outage probabilities of the weak and strong users, respectively.

Assuming independent decoding events for the paired users, the successful transmission probability of one NOMA pair is

\begin{equation}
P_{success}
=
(1-P_{out,w})
(1-P_{out,s}).
\end{equation}

Accordingly, the outage probability of the NOMA pair is

\begin{equation}
P_{pair}
=
1-
(1-P_{out,w})
(1-P_{out,s}).
\label{eq:pairout}
\end{equation}

Equation (\ref{eq:pairout}) indicates that the outage probability decreases whenever either user's decoding reliability improves.

For the proposed Hybrid SIC/CIC framework, the pair outage probability automatically adapts to the cache status of both users through the unified SINR expressions. Consequently,

\begin{itemize}

\item when neither user possesses cached information ($c_w=0$, $c_s=0$), (\ref{eq:pairout}) reduces to the conventional SIC-NOMA pair outage;

\item when only the weak user performs cache-aided interference cancellation ($c_w=1$, $c_s=0$), the weak-user outage decreases while the strong-user outage remains unchanged;

\item when only the strong user performs cache-aided interference cancellation ($c_w=0$, $c_s=1$), the strong-user outage decreases;

\item when both users exploit locally cached contents ($c_w=1$, $c_s=1$), the outage probabilities of both users are minimized, resulting in the lowest pair outage probability.

\end{itemize}

Therefore, the proposed Hybrid SIC/CIC receiver improves the reliability of every NOMA pair by dynamically selecting the appropriate interference cancellation mechanism according to the instantaneous cache availability.
```
```latex
%====================================================
\subsection{End-to-End System Outage Probability}
%====================================================

While the pair outage probability evaluates the reliability of an individual NOMA pair, the overall network performance depends on the successful transmission of all scheduled NOMA pairs. Therefore, the end-to-end outage probability is adopted as the primary reliability metric of the proposed Hybrid SIC/CIC framework.

Assume that the network contains $L$ independent NOMA pairs. Let
$P_{pair,i}$ denote the outage probability of the $i$-th pair.

The successful transmission probability of the entire system is

\begin{equation}
P_{success}^{sys}
=
\prod_{i=1}^{L}
\left(1-P_{pair,i}\right).
\label{eq:success_system}
\end{equation}

Consequently, the end-to-end outage probability is obtained as

\begin{equation}
P_{E2E}
=
1-
\prod_{i=1}^{L}
\left(1-P_{pair,i}\right).
\label{eq:e2e}
\end{equation}

For a sufficiently large number of user pairs, the average system outage can alternatively be expressed as

\begin{equation}
P_{sys}
=
\frac{1}{L}
\sum_{i=1}^{L}
P_{pair,i}.
\label{eq:avg_system}
\end{equation}

Equation (\ref{eq:e2e}) demonstrates that the overall network reliability depends on the outage performance of every scheduled NOMA pair. Since the proposed Hybrid SIC/CIC receiver reduces the outage probability of both weak and strong users through cache-assisted interference cancellation, the end-to-end outage probability is significantly lower than that of the conventional SIC-based NOMA system.

%====================================================
\subsection{Dynamic Cache-Aware Power Allocation}
%====================================================

Unlike conventional NOMA systems that employ fixed power allocation coefficients, the proposed framework dynamically adjusts the transmit power according to the instantaneous channel conditions and cache availability of each NOMA pair.

Let

\begin{equation}
\alpha_{w,m}+\alpha_{s,m}=1,
\label{eq:power_sum}
\end{equation}

where

\begin{equation}
0<\alpha_{s,m}<\alpha_{w,m}<1.
\end{equation}

The power allocation coefficients are determined according to the receiver operating mode.

\textbf{Mode 1: Conventional SIC}

When no cache-assisted interference cancellation is available,

\[
c_w=0,\qquad c_s=0,
\]

the weak user requires higher transmit power to overcome multi-user interference. Therefore,

\begin{equation}
\alpha_{w,m}\gg\alpha_{s,m}.
\end{equation}

This corresponds to the conventional NOMA power allocation strategy.

\textbf{Mode 2: Weak-User CIC}

If

\[
c_w=1,
\]

the weak user removes the interference using its cached information. Hence, the minimum power required to satisfy the target SINR becomes

\begin{equation}
\alpha_{w,m}^{min}
=
\frac{\gamma_{th}\sigma^2}
{P_m|h_{w,m}|^2}.
\label{eq:weakpower}
\end{equation}

The remaining transmit power can therefore be allocated to the strong user, improving the overall system throughput.

\textbf{Mode 3: Strong-User CIC}

When

\[
c_s=1,
\]

the residual SIC interference is completely eliminated. Consequently, the feasible power region expands as

\begin{equation}
\alpha_{w,m}^{max}
=
1-
\frac{\gamma_{th}\sigma^2}
{P_m|h_{s,m}|^2}.
\label{eq:strongpower}
\end{equation}

This additional degree of freedom allows more efficient power sharing between the paired users.

\textbf{Mode 4: Full Hybrid CIC}

When both users possess the interfering content,

\[
c_w=c_s=1,
\]

the proposed algorithm jointly determines the optimal power coefficients

\[
(\alpha_w,\alpha_s),
\]

such that the outage probability is minimized while maintaining the required Quality-of-Service constraints for both users.

Consequently, the proposed cache-aware power allocation dynamically adapts the transmit power according to both channel quality and cache status, leading to improved spectral efficiency, lower outage probability, and higher energy efficiency.

%====================================================
\subsection{Optimization Problem}
%====================================================

The objective of the proposed Hybrid SIC/CIC framework is to jointly optimize the power allocation coefficients for all NOMA pairs while exploiting cache-assisted interference cancellation to improve network reliability.

The optimization problem is formulated as

\begin{equation}
\begin{aligned}
\max_{\{\alpha_{w,m},\alpha_{s,m}\}}
\quad
&
\sum_{m=1}^{M}
\left(
R_{w,m}+R_{s,m}
\right)
\\
\text{s.t.}
\quad
&
P_{out,w}\le\epsilon_w,
\\
&
P_{out,s}\le\epsilon_s,
\\
&
R_{w,m}\ge R_{th},
\\
&
R_{s,m}\ge R_{th},
\\
&
\alpha_{w,m}+\alpha_{s,m}=1,
\\
&
0<\alpha_{s,m}<\alpha_{w,m}<1,
\\
&
P_m\le P_{max},
\\
&
c_w,c_s\in\{0,1\},
\end{aligned}
\label{eq:optimization}
\end{equation}

where

\begin{itemize}

\item $R_{th}$ denotes the minimum QoS data rate,

\item $\epsilon_w$ and $\epsilon_s$ represent the maximum allowable outage probabilities,

\item $P_{max}$ is the BS transmit power budget,

\item $c_w$ and $c_s$ indicate the cache status of the paired users.

\end{itemize}

The optimization problem in (\ref{eq:optimization}) is non-convex because the achievable rates and outage probabilities are nonlinear functions of the power allocation coefficients and the cache-assisted interference cancellation states. Therefore, instead of relying on computationally expensive iterative optimization methods, this work employs a Deep Q-Network (DQN)-based learning framework to determine near-optimal power allocation decisions. The DQN agent continuously observes the network state, including channel gains, cache availability, and user QoS requirements, and learns an adaptive power allocation policy that maximizes the long-term system sum rate while simultaneously minimizing the end-to-end outage probability under the given transmit power constraints.
```


%====================================================
\subsection{Achievable Rate}
%====================================================

Let

\begin{equation}
\rho=\frac{P_m}{B N_0}
\end{equation}

denote the transmit SNR.

The achievable rate of the weak user is

\begin{equation}
R_{w,m}
=
B\log_2
\left(
1+\gamma_{w,m}
\right),
\end{equation}

while the strong user's rate is

\begin{equation}
R_{s,m}
=
B\log_2
\left(
1+\gamma_{s,m}
\right).
\end{equation}

The total NOMA throughput is

\begin{equation}
R_{\mathrm{sum}}
=
\sum_{m=1}^{M}
\left(
R_{w,m}+R_{s,m}
\right).
\end{equation}

%====================================================
\subsection{Resource Block Allocation}
%====================================================

Let

\begin{equation}
\beta_{k,m}
=
\begin{cases}
1,& \text{if RB }m\text{ is assigned to user }k,\\
0,& \text{otherwise}.
\end{cases}
\end{equation}

Each RB can serve at most two NOMA users:

\begin{equation}
\sum_{k=1}^{K}
\beta_{k,m}
\le 2,
\qquad \forall m.
\end{equation}

The total power allocation satisfies

\begin{equation}
\sum_{k=1}^{K}
\sum_{m=1}^{M}
\alpha_{k,m}
\le 1.
\end{equation}


% IV. PROPOSED DRL CACHING FRAMEWORK
%=============================================================
\section{Proposed DRL Caching Framework}

\subsection{Proposed Algorithm Initialization}
To effectively outline the online placement boundaries applied immediately following the Hybrid Cache-Aided TDMA-NOMA transmission layer outlined in the System Model, the formalized execution parameters are mapped strictly within Algorithm \ref{alg:dqn}.

\begin{algorithm}
\caption{Hybrid Cache-Aware TDMA-NOMA Double Dueling DQN Cache Training}
\label{alg:dqn}
\begin{algorithmic}[1]
\REQUIRE User set $K=200$, File catalog $N=2000$, Cache capacity $C=200$, Maximum training episodes $E = 500$
\STATE \textbf{Variables:} Online network $\theta$ (the primary model evaluating real-time actions), Target network $\theta^-$ (delayed clone for stability), Replay buffer $\mathcal{B}$
\STATE Initialize online network $\theta$, target network $\theta^- \leftarrow \theta$, and bounded replay buffer $\mathcal{B}$
\STATE Assign exploration factor $\epsilon \leftarrow 1.0$, global timeline $t \leftarrow 0$
\FOR{episode $e = 1$ to $E = 500$}
    \STATE Reset cache buffer; instantiate fresh user positions, channel conditions, and cache states
    \STATE Generate random request stream $(u_i, f_i)_{i=1}^{T}$ for $T=10000$
    \FOR{each dynamic request $(u_i, f_i)$}
        \STATE Extract feature state $s_t \in \mathbb{R}^{14}$
        \STATE With probability $\epsilon$: $a_t \leftarrow$ Random Action; else $a_t \leftarrow \arg\max_a Q(s_t, a; \theta)$
        \IF{$a_t = 1$ and target file absent from cache}
            \STATE Insert file $f_i$; eject least popular entry if memory full
        \ENDIF
        \STATE Sort active users according to channel gain
        \STATE Form strong--weak user pairs using channel disparity and cache availability
        \STATE Allocate each selected pair to the corresponding TDMA transmission slot
        \STATE Perform Hybrid TDMA-NOMA superposition transmission
        \STATE Apply Cache-Aided Interference Cancellation (CIC) followed by SIC whenever required
        \STATE Compile cache-hit ratio, CIC success, achievable sum-rate and energy-efficiency metrics; observe reward $r_t$ via \eqref{eq:reward}
        \STATE Queue transition $(s_t, a_t, r_t, s_{t+1})$ into $\mathcal{B}$ with max priority flag
        \IF{$t > t_{\text{warmup}}$ and $(t \bmod 4 = 0)$}
            \STATE Extract IS-weighted transition minibatch from $\mathcal{B}$
            \STATE Calculate Double-Q target $y_t$ leveraging \eqref{eq:ddqn}
            \STATE Update bounds $\theta$ via Huber Loss backpropagation
            \STATE Soft-blend target network $\theta^-$ using $\tau = 0.005$
        \ENDIF
        \STATE Decay $\epsilon$ linearly; advance global tick $t \leftarrow t + 1$
    \ENDFOR
\ENDFOR
\end{algorithmic}
\end{algorithm}

\subsection{MDP Formulation}
We formulate the online cache placement problem as a continuous Markov Decision Process (MDP) defined by the tuple $\langle \mathcal{S}, \mathcal{A}, \mathcal{R}, \mathcal{P}, \gamma \rangle$.

\textbf{State space:} Upon receiving a request at step $t$, the extracted state vector $s_t \in \mathbb{R}^{14}$ encodes:
(i) Normalized popularity rank of the requested file,
(ii) Localized recent request frequency scalar of the target file,
(iii) Binary cache membership boolean indicator,
(iv) Cache storage occupancy ratio,
(v--vii) Scalar mean, min, and max historical popularity traits of currently cached files,
(viii) Binary quality-improvement boolean indicator for the present cache swap,
(ix) Popularity ratio relative to the absolute least popular cached file,
(x) Moving-average CIC potential rate from recent NOMA pairing history,
(xi--xii) Mean and standard deviation distributions of real-time channel gains,
(xiii) Historical NOMA transmission success tracking rate, and
(xiv) Evaluated file request diversity ratio bounded within the moving window.

\textbf{Action space:} In contrast to traditional policies evaluating $\mathcal{O}(N)$ eviction slots, our action definition is strictly binary: $a_t \in \{0, 1\}$. $a_t = 1$ forces the cache to encode the newly requested file (triggering an internal heuristic lowest-popularity eviction if full), whereas $a_t = 0$ actively skips the caching stage. This reduces learning complexity to $\mathcal{O}(1)$.

\textbf{Reward function:} The immediate state-action transition reward $r_t$ isolates strict decision boundary qualities:

\begin{equation}
r_t =
\begin{cases}
+2.0 & \text{cache hit (correct prior decision)} \\
+1.5 & \text{cache top-10\% popular file} \\
+0.8 & \text{cache top-30\% popular file} \\
+0.2 & \text{cache top-50\% popular file} \\
-0.5 & \text{cache bottom-50\% (unpopular file)} \\
+1.5 & \textbf{CIC enabled by this placement} \\
-0.5 & \text{skip a popular file with free space} \\
+0.1 & \text{correctly skip unpopular file}
\end{cases}
\label{eq:reward}
\end{equation}

Furthermore, explicit eviction quality bonuses ($\pm 0.5$) are applied exclusively when the storage bounds actuate a deletion (i.e., verifying if the swapped file historically underperformed the insertion candidate).

\subsection{Double Dueling DQN Architecture}
The central Q-function estimator relies on a \textit{Dueling DQN} topology \cite{b_wang_dueling} which intelligently segregates the raw state value $V(s)$ from action-specific advantages $A(s, a)$:

\begin{equation}
Q(s, a; \theta) = V(s; \theta_V) + \left[ A(s, a; \theta_A) - \frac{1}{|\mathcal{A}|}\sum_{a'} A(s, a'; \theta_A) \right]
\label{eq:dueling}
\end{equation}

The structural network dictates: Input($14$) $\rightarrow$ FC($64$)-ReLU $\rightarrow$ FC($32$)-ReLU. This trunk subsequently bifurcates into an autonomous value stream FC($16$)-ReLU-FC($1$) and an equivalent advantage stream FC($16$)-ReLU-FC($2$). To mitigate Q-value overestimation phenomena, the \textit{Double DQN} paradigm \cite{b_van_hasselt} decouples target action selection from mathematical evaluation:

\begin{equation}
y_t = r_t + \gamma \cdot Q_{\text{target}}\!\left(s_{t+1}, \arg\max_{a} Q_{\text{online}}(s_{t+1}, a)\right)
\label{eq:ddqn}
\end{equation}
enforcing an infinite horizon discount factor $\gamma = 0.99$.

\subsection{Prioritized Experience Replay}
State transitions are pushed into an asynchronous prioritized replay buffer \cite{b_schaul_per} comprising $50,000$ slots with a sampling exponent $\alpha_{\text{PER}} = 0.6$. Importance sampling (IS) weights equipped with $\beta$ annealed steadily from $0.4$ to $1.0$ over $5\times 10^6$ training steps adjust for non-uniform loss updates.

\subsection{Computational Complexity}
A critical advantage of the proposed framework is its runtime efficiency. While optimal proactive caching requires an NP-hard combinatorial search over $\mathcal{O}\binom{N}{C}$ subsets, the trained DRL agent operates in constant continuous time. The forward pass of the highly compact $14 \rightarrow 64 \rightarrow 32 \rightarrow 16 \rightarrow 2$ network involves merely $\approx 3,700$ FLOPs per request, achieving $\mathcal{O}(1)$ execution complexity. This ensures strict adherence to the ultra-low latency constraints demanded by 6G URLLC (Ultra-Reliable Low-Latency Communications) standards.


%=============================================================
% V. SIMULATION SETUP
%===========================================================
\section{Simulation Setup}

The performance of the proposed Hybrid SIC/CIC-enabled DQN framework is evaluated through Monte Carlo simulations implemented in Python. Unless otherwise stated, the simulation parameters summarized in Table~\ref{tab:params} are adopted throughout the paper. The selected values follow commonly used assumptions in cache-aided NOMA literature and are consistent with \cite{Dani2022, Ding2017, Liu2018}.

\begin{table}[!t]
\caption{Simulation Parameters}
\label{tab:params}
\centering
\begin{tabular}{ll}
\hline
\textbf{Parameter} & \textbf{Value} \\
\hline
Cell topology & Single-cell downlink \\
Number of users ($K$) & 200 \\
Cell radius ($R$) & 500 m \\
Number of Resource Blocks ($M$) & 100 \\
Users per RB & 2 (NOMA pair) \\
Content library size ($N$) & 2000 files \\
Cache size per user ($C$) & 200 files (10\%) \\
Content popularity & Zipf distribution \\
Zipf exponent ($\alpha$) & 1.0 \\
Channel model & Rayleigh fading \\
Path-loss exponent ($\beta$) & 3.5 \\
Noise power ($\sigma^2$) & $10^{-9}$ W \\
Maximum BS transmit power ($P_{\max}$) & 2 W (33 dBm) \\
Target rate ($R_{\mathrm{th}}$) & 0.3 bps/Hz \\
SINR threshold & $2^{R_{\mathrm{th}}}-1$ \\
Residual SIC factor ($\zeta$) & 0.05 \\
Initial power coefficients & $(\alpha_w,\alpha_s)=(0.8,0.2)$ \\
User pairing & Extreme pairing \\
Interference cancellation & Hybrid SIC/CIC \\
Power allocation & DQN-based adaptive \\
\hline
\multicolumn{2}{c}{\textbf{DQN Hyperparameters}}\\
\hline
Learning algorithm & Double Dueling DQN \\
Discount factor ($\gamma$) & 0.99 \\
Learning rate & $10^{-3}$ \\
Optimizer & Adam \\
Replay memory & 50,000 \\
Mini-batch size & 64 \\
Target network update & Every 200 steps \\
Initial exploration ($\epsilon$) & 1.0 \\
Final exploration ($\epsilon_{\min}$) & 0.01 \\
Exploration decay & Linear \\
Training episodes & 500 \\
\hline
\multicolumn{2}{c}{\textbf{Performance Evaluation}}\\
\hline
Monte Carlo realizations & $10^{3}$ per SNR \\
SNR range & $-10$ to $30$ dB \\
Performance metrics &
Sum rate, Outage probability,\\
& Cache hit ratio, Energy efficiency \\
\hline
\end{tabular}
\end{table}

The proposed Hybrid SIC/CIC-enabled Double Dueling DQN framework is compared with four benchmark schemes:

\begin{itemize}

\item \textbf{No Cache-NOMA}: Conventional power-domain NOMA employing imperfect SIC without caching.

\item \textbf{LRU Cache}: Cache-enabled NOMA using the Least Recently Used replacement policy.

\item \textbf{LFU Cache}: Cache-enabled NOMA using the Least Frequently Used replacement policy.

\item \textbf{Top-$K$ Cache}: Popularity-aware caching that stores the most frequently requested files according to the Zipf distribution.

\end{itemize}

The proposed method dynamically learns cache-aware power allocation policies using the Hybrid SIC/CIC receiver, and its performance is evaluated in terms of achievable sum rate, outage probability, cache hit ratio, and energy efficiency.
%=============================================================
% VI. RESULTS AND DISCUSSION
%=============================================================
\section{Results and Discussion}

\subsection{DQN Training Convergence Profiles}

Fig.~\ref{fig:training} visualizes the DQN stabilization properties continuously polled across 500 isolated episodes. The fundamental cache hit rate traces a robust ascent from $0.670$ strictly plateauing near $0.714$ by episode 150. Isolated test phases remain uniformly 1--2\% elevated above the training loss curves, proving definitive robustness against overfitting.

\begin{figure}[htbp]
\centerline{\includegraphics[width=\columnwidth]{dqn_training_curves.png}}
\caption{DQN offline training bounds mapped over 500 episodes (encompassing 5M simulation steps). Features cache hit rate metrics, $\epsilon$-decay constraints, and cumulative Huber loss gradients.}
\label{fig:training}
\end{figure}

The empirical loss curve maps the expected architectural double-spike pattern characteristic of PER-enhanced Dueling networks: instantaneous losses scale to $\approx 1200$ near episode 80 as exploration phases contract into strict exploitation, subsequently crumbling into near-zero residuals terminating around episode 150. Consequently, the modeled CIC benefit tracking line exhibits an identical stabilization lock ascending linearly from $\approx 3.8$ into a sustained $4.1$. We observe the spectral boundary capping gracefully near an ideal 28 bps/Hz threshold.

\subsection{Policy Verification Framework (Monte Carlo Analysis)}

The discrete Monte Carlo comparative boundaries are isolated in Fig.~\ref{fig:comparison} and thoroughly summarized in Table~\ref{tab:comparison}.

\begin{figure}[htbp]
\centerline{\includegraphics[width=\columnwidth]{comparison_plots.png}}
\caption{Rigorous box-plot variance evaluations mapping the proposed DQN protocol systematically against Top-K, LFU, LRU, and strictly randomized caching baselines over 100 iterations.}
\label{fig:comparison}
\end{figure}

\begin{table}[htbp]
\centering
\caption{Monte Carlo Mean $\pm$ Standard Deviation Matrix}
\label{tab:comparison}
\begin{tabular}{@{}lcccc@{}}
\toprule
\textbf{Policy} & \textbf{Hit Rate} & \textbf{Outage} & \textbf{CIC Rate} & \textbf{SE (bps/Hz)} \\
\midrule
\textbf{DQN} & \textbf{0.7038$\pm$0.004} & 0.0649$\pm$0.018 & \textbf{3.776$\pm$0.091} & 26.09$\pm$0.51 \\
TopK & 0.7092$\pm$0.005 & 0.0648$\pm$0.018 & 0.715$\pm$0.035 & \textbf{26.72$\pm$0.62} \\
LFU & 0.6847$\pm$0.005 & 0.0649$\pm$0.018 & 0.686$\pm$0.037 & 24.04$\pm$0.53 \\
LRU & 0.6122$\pm$0.005 & 0.0645$\pm$0.018 & 0.627$\pm$0.034 & 18.11$\pm$0.40 \\
Random & 0.5598$\pm$0.006 & 0.0644$\pm$0.018 & 0.575$\pm$0.035 & 15.05$\pm$0.34 \\
\bottomrule
\end{tabular}
\end{table}

\textbf{Cache Hit Rate Insights:} The DRL abstraction maintains an impressive absolute hit rate of 70.38\%, trailing the impossible theoretical `TopK-Zipf` oracle threshold (70.92\%) by a negligible 0.54\%. Contextually, this represents a massive 26\% performance inflation over LRU standards (61.22\%) achieved autonomously by the agent without any pre-configured mathematical knowledge defining the Zipf payload architecture.

\textbf{CIC Efficiency Dominance:} Evaluating NOMA interference boundaries, the reinforcement framework maps an average CIC exploitation rate scaling uniformly to 3.776, definitively outclassing the static Top-K boundary (0.715). This represents an irrefutable \textbf{5.28x system improvement}, solidifying the underlying premise that exclusively assigning rewards mapped conditionally to multi-user NOMA cancellation phenomena forces behavioral traits impossible to emulate in strict mathematical placement algorithms.

\subsection{Cache-Aided Environment Metrics: Detailed SNR Analysis}
To provide a comprehensive evaluation across varying signal conditions, the multi-dimensional tracking profile is decomposed into explicit NOMA performance thresholds. The following 9 analytical subsections (Fig. \ref{fig:sum_rate} - Fig. \ref{fig:energy}) correspond to the isolated evaluation variables plotted across the discrete $-10$ dB to $+30$ dB SNR grid.

\begin{figure}[htbp]
\centerline{\includegraphics[width=\columnwidth]{fig_sum_rate.png}}
\caption{System Sum-Rate vs. SNR (dB), demonstrating proactive caching throughput expansions.}
\label{fig:sum_rate}
\end{figure}
\subsubsection{Peak Capacity Yields (Sum-Rate)}
Maximized under a localized signal envelope representing SNR = 30 dB, the Dueling module pushes maximum channel throughputs hitting 18.17 bps/Hz, trailing the static oracle limitation (18.18 bps/Hz) by fractions of a boundary. By contrast, stripped identically of edge deployment caches (No-Cache boundary), general throughputs plateau firmly near 5.59 bps/Hz. Overall, proactive deep placements attribute massive multi-user capability amplifications directly correlating to \textbf{225\% throughput ceiling explosions}, validating theoretical capacity limits discussed in \cite{b_cao_joint}.

\begin{figure}[htbp]
\centerline{\includegraphics[width=\columnwidth]{fig_outage.png}}
\caption{Logarithmic Outage Probability vs. SNR (dB), highlighting transmission reliability.}
\label{fig:outage}
\end{figure}
\subsubsection{Network Outage Collapse}
Logarithmic outage traces visually collapse the baseline No-Cache probability ceilings from $0.085$ firmly flatlining around near-optimal boundaries measuring $0.006$. Stated directly, proactively leveraging NOMA coupled intimately with cached contents isolates an overall connectivity robustness amplification factor measuring strictly \textbf{14x lower drop-outs} at localized extreme signal bands. This confirms the outage minimization models established by Zhang \textit{et al.} \cite{b_zhang_stochastic}.

\begin{figure}[htbp]
\centerline{\includegraphics[width=\columnwidth]{fig_hit_rate.png}}
\caption{Base-Station Cache Hit Rate reliability across varying multi-path channel profiles.}
\label{fig:hit_rate}
\end{figure}
\subsubsection{Cache Hit Rate Consistency}
The isolated cache hit probability curve demonstrates the DQL agent reliably securing $\approx 70.38\%$ predictive success regardless of the underlying transmission SNR geometry. While dynamic baselines like LRU degrade significantly (61.2\%) due to heavy temporal bias, the optimized deep feature selection successfully tracks the explicit Zipf distribution, bounding directly below the perfect subset Top-K limit \cite{b_breslau_zipf}.

\begin{figure}[htbp]
\centerline{\includegraphics[width=\columnwidth]{fig_weak_rate.png}}
\caption{Cell-edge (Weak) User Rate expansions driven perfectly by CIC application.}
\label{fig:weak_rate}
\end{figure}
\subsubsection{Weak User (Cell-Edge) Rate Expansion}
A pervasive constraint driving standard NOMA deployment limits involves punishing weak-user capacity floors under severe path interference. By predicting weak user demands and injecting Cache-Aided Interference Cancellation (CIC), weak user data rates aggressively vault far beyond traditional power-domain allocations, equalizing service bounds structurally \cite{b_survey_bepari}.

\begin{figure}[htbp]
\centerline{\includegraphics[width=\columnwidth]{fig_strong_rate.png}}
\caption{Cell-center (Strong) User Rate mapping incorporating imperfect SIC overheads.}
\label{fig:strong_rate}
\end{figure}
\subsubsection{Strong User (Cell-Center) Delivery}
Simultaneously, the center-placed strong user benefits inherently from the identical NOMA coupling structure. Reconstructing the weak-user multi-dimensional interference dynamically from localized cache memory heavily limits imperfect SIC error chains. This generates throughput slopes mimicking optimal point-to-point orthogonality models, bypassing interference barriers \cite{b_dai_noma}.

\begin{figure}[htbp]
\centerline{\includegraphics[width=\columnwidth]{fig_ber_weak.png}}
\caption{Bit Error Rate (BER) mitigation mapping unified decoding failure reduction for Weak Users.}
\label{fig:ber_weak}
\end{figure}
\subsubsection{Bit Error Rate (BER) Suppression - Weak User}
Bit Error Rate trajectories evaluated specifically for the weak user confirm a strict 1--2 order of magnitude collapse in hardware-level decoding fatalities specifically traversing the $0$ to $15$ dB boundary envelope. Relieving standard Successive Interference Cancellation hardware limits actively bypasses fatal cascading bit-errors \cite{b_li_drl}.

\begin{figure}[htbp]
\centerline{\includegraphics[width=\columnwidth]{fig_ber_strong.png}}
\caption{Bit Error Rate (BER) trajectories for Strong Users leveraging perfect CIC bounds.}
\label{fig:ber_strong}
\end{figure}
\subsubsection{Bit Error Rate (BER) Suppression - Strong User}
Similarly, the center-placed strong user bounds approach mathematical minimum interference ceilings. By converting cache hits directly into flawless multi-user cancellation, the strong user eliminates the foundational 5\% imperfect SIC tolerance ($\zeta = 0.05$) which strictly bounds baseline random models, equating edge caching capabilities to dynamic hardware upgrades \cite{b_survey_bepari}.

\begin{figure}[htbp]
\centerline{\includegraphics[width=\columnwidth]{fig_fairness.png}}
\caption{Jain's Fairness Index plotted over SNR, measuring internal resource equalization.}
\label{fig:fairness}
\end{figure}
\subsubsection{System Fairness Equalization}
Jain's Fairness Index strictly quantifies internal resource equality constraints across participating paired spectrums. While baseline transmission models rapidly diverge below index levels of $0.62$, the CIC-incentivized cache perfectly equalizes power-sharing bounds returning a peak theoretical fairness alignment hitting $0.957$. 

\begin{figure}[htbp]
\centerline{\includegraphics[width=\columnwidth]{fig_energy.png}}
\caption{Unified Network Energy Efficiency mapping calculated in bits per Joule (bits/J).}
\label{fig:energy}
\end{figure}
\subsubsection{Network Energy Efficiency}
Consequently, localized cached payload deliveries explicitly bypass base-station retransmission cycling and backhaul communication chains entirely, spiking localized energy boundaries. The intelligent framework extracts a staggering peak evaluation of $18.66$ bits per Joule (bits/J) natively surpassing purely dynamic LRU routing policies while outperforming non-cached models geometrically \cite{b_cao_joint}.

%=============================================================
% VII. CONCLUSION
%=============================================================
\section{Conclusion}

This paper has definitively formalized a Double Dueling Deep Q-Network pipeline structured specifically around proactively generating localized caching policies mapped mathematically into NOMA network transmission paths for dynamic 6G infrastructure. Replacing infinite algorithmic decision boundaries with a strict binary resolution model interacting sequentially with a heavily condensed 14-dimensional channel state, the module autonomously isolates optimal delivery sequences prioritizing the exclusive reconstruction elements driving Cache-Aided Interference Cancellation (CIC). The proposed hybrid SIC/CIC framework enables both weak-user and strong-user interference cancellation modes, while the cache-aware dynamic power allocation strategy adapts transmit power coefficients in real-time to exploit the altered interference structure---delivering additional energy savings and throughput gains beyond what static power allocation can achieve. Independent Monte Carlo validations effectively verify sustained performance evaluations matching explicit ideal oracle limitations down to 0.54\% hit probability variances. Simultaneously isolating multi-user interplay rewards results directly inside unprecedented 5.28x factor CIC expansions. These properties effectively guarantee comprehensive reductions dictating absolute operational downtime (reducing outage thresholds uniformly 14x against base limitations), firmly justifying Deep RL frameworks acting centrally inside interconnected 6G edge communication layers.

Future exploration trajectories anticipate formulating scaled reinforcement interactions targeting dynamic user mobility topologies combined intimately via localized federated network coordination deployments spanning wide-area multi-cell coverage bounds.


%=============================================================
% ACKNOWLEDGMENT
%=============================================================
\section*{Acknowledgment}
The authors respectfully acknowledge the critical project steering and invaluable supervisory oversight uniquely contributed by Dr. Biswajit Ghosh representing the Department of Information Technology situated within Techno Main Salt Lake, directly facilitating the core developments powering this Project.


%=============================================================
% REFERENCES
%=============================================================
\bibliographystyle{IEEEtran}
%\bibliography{ref}

\end{document}